\documentclass[12pt]{article}
\usepackage[utf8]{inputenc}
\usepackage[T1]{fontenc}
\usepackage[margin=1in]{geometry}
\usepackage{setspace}
\usepackage{xcolor}
\usepackage[colorlinks=true, urlcolor=blue!60!black, linkcolor=blue!60!black]{hyperref}
\setlength{\parindent}{2em}
\setlength{\parskip}{1.25em}
\renewcommand{\baselinestretch}{1.0}
\title{\textit{Trump v. Hawaii}, 585 U.S. 667 (2018)}
\author{Case Brief}
\date{}
\begin{document}
\maketitle

\section{Case Metadata}

\begin{itemize}
  \item \textbf{Case Name:} \textit{Trump v. Hawaii}
  \item \textbf{Citation:} \href{https://supreme.justia.com/cases/federal/us/585/17-965/}{585 U.S. 667 (2018)}
  \item \textbf{Court and Dates:} Supreme Court of the United States; argued April 25, 2018; decided June 26, 2018
  \item \textbf{Chapter:} Chapter 2: Separation of Powers and the Executive
  \item \textbf{Vote:} 5--4 (Majority: Roberts, C.J., joined by Kennedy, Thomas, Alito, and Gorsuch, JJ.; Dissenting: Breyer, J., joined by Kagan, J.; Sotomayor, J., joined by Ginsburg, J.)
  \item \textbf{Author of the Opinion of the Court:} Chief Justice John G. Roberts, Jr.
  \item \textbf{Concurring Opinions:} Justice Anthony M. Kennedy; Justice Clarence Thomas
  \item \textbf{Dissenting Opinions:} Justice Stephen G. Breyer (joined by Justice Elena Kagan); Justice Sonia Sotomayor (joined by Justice Ruth Bader Ginsburg)
  \item \textbf{Constitutional and Statutory Provisions:} U.S. Const. amend. I (Establishment Clause); Immigration and Nationality Act (INA) \S~212(f), 8 U.S.C. \S~1182(f); INA \S~202(a)(1)(A), 8 U.S.C. \S~1152(a)(1)(A)
  \item \textbf{Primary Source Links:} \href{https://www.supremecourt.gov/opinions/17pdf/585us2r70_jgkn.pdf}{SupremeCourt.gov Official Proof}; \href{https://supreme.justia.com/cases/federal/us/585/17-965/}{Justia}; \href{https://www.law.cornell.edu/supremecourt/text/17-965}{Cornell LII}
\end{itemize}

\section{Facts and Procedural History}

Shortly after taking office in January 2017, President Donald J. Trump issued Executive Order No. 13769 (``EO-1'') and later Executive Order No. 13780 (``EO-2''), temporarily restricting entry from several predominantly Muslim nations while directing the Department of Homeland Security (DHS) to conduct a worldwide review of foreign governments' identity-management and information-sharing practices. Following that multi-agency review, President Trump issued Presidential Proclamation No. 9645 on September 24, 2017, invoking INA \S~212(f), 8 U.S.C. \S~1182(f), which empowers the President to ``suspend the entry of all aliens or any class of aliens'' whenever he ``finds'' that their entry ``would be detrimental to the interests of the United States.''

Proclamation No. 9645 imposed tailored entry restrictions on nationals of eight countries whose vetting protocols DHS deemed deficient or that posed terrorism risks: Chad, Iran, Libya, North Korea, Somalia, Syria, Venezuela, and Yemen. The Proclamation exempted lawful permanent residents and asylees, established a case-by-case waiver program for undue hardship, and mandated DHS reviews every 180 days. Following the first review, President Trump lifted restrictions on Chad in April 2018 (Proclamation No. 9723), leaving seven covered countries---five Muslim-majority nations plus North Korea and Venezuela.

The State of Hawaii, three U.S. citizens with affected foreign relatives (including Dr. Ismail Elshikh), and the Muslim Association of Hawaii sued in federal district court, alleging that the Proclamation exceeded the President's authority under \S~1182(f), violated the INA's ban on nationality discrimination in immigrant-visa issuance (8 U.S.C. \S~1152(a)(1)(A)), and violated the Establishment Clause. To support the constitutional claim, plaintiffs cited President Trump's campaign and post-inauguration statements calling for a ``total and complete shutdown of Muslims entering the United States'' and characterizing Islam in hostile terms. The District Court entered a nationwide preliminary injunction, and the Ninth Circuit affirmed on statutory grounds without reaching the Establishment Clause (878 F.3d 662). After staying the injunction pending appeal, the Supreme Court granted certiorari.

\section{Issue(s) Presented}

\begin{enumerate}
  \item \textbf{Statutory Authority:} Did Presidential Proclamation No. 9645 exceed the President's delegated authority under 8 U.S.C. \S~1182(f) or conflict with the nondiscrimination mandate of 8 U.S.C. \S~1152(a)(1)(A)?
  \item \textbf{Establishment Clause:} What standard of judicial review governs an Establishment Clause challenge to a facially neutral presidential proclamation restricting the entry of foreign nationals, and did plaintiffs demonstrate a likelihood of success on the merits given the President's extrinsic statements?
\end{enumerate}

\section{Holding}

\textbf{Reversed and remanded (5--4).} Chief Justice Roberts held that: (1) Proclamation No. 9645 was a lawful exercise of the President's broad statutory discretion under 8 U.S.C. \S~1182(f) and did not violate \S~1152(a)(1)(A); (2) under deferential review derived from \textit{Kleindienst v. Mandel}, 408 U.S. 753 (1972), and rational-basis scrutiny, the Proclamation did not violate the Establishment Clause because it was plausibly grounded in a legitimate national-security rationale independent of religious animus; and (3) \textit{Korematsu v. United States}, 323 U.S. 214 (1944), is expressly repudiated as ``gravely wrong the day it was decided.''

\section{Reasoning}

\textbf{1. Broad Delegation Under \S~1182(f).} Assuming without deciding that plaintiffs' statutory claims were justiciable, Roberts concluded that \S~1182(f)'s plain text ``exudes deference to the President in every clause.'' The statute's sole prerequisite---that the President ``find'' entry would be detrimental to U.S. interests---was amply satisfied by the Proclamation's 12-page explanation of the worldwide multi-agency vetting review. Rejecting plaintiffs' argument that \S~1182(f) permits only temporary emergency measures, the Court noted that past Presidents (including Reagan and Clinton) had used the statute to impose nationality-based entry suspensions without fixed end dates.

\textbf{2. Harmonizing \S~1182(f) and \S~1152(a)(1)(A).} Section 1152(a)(1)(A) prohibits nationality discrimination ``in the issuance of an immigrant visa.'' Roberts distinguished \textit{entry} eligibility (governed by \S~1182) from \textit{visa issuance} (governed by \S~1152): once a foreign national is deemed inadmissible under \S~1182(f), he is ineligible for a visa as a threshold matter, leaving no conflict between the two sections.

\textbf{3. Establishment Clause and Deferential Rational-Basis Review.} Assuming plaintiffs had standing based on separation from foreign relatives, the Court declined to apply domestic Establishment Clause tests that strike down policies whenever extrinsic statements suggest religious animus. Because foreign nationals abroad lack constitutional entry rights and immigration admissions implicate core national-security and foreign-affairs judgments of the political branches, \textit{Mandel} normally limits courts to verifying a ``facially legitimate and bona fide reason.'' Even looking behind the face of the Proclamation at the President's statements under rational-basis review, the Court upheld the policy if it could ``reasonably be understood to result from a justification independent of unconstitutional grounds.'' Here, the Proclamation was textually neutral, rested on an extensive Cabinet-level vetting review, removed three Muslim-majority countries (Iraq, Sudan, and Chad) over time, exempted most of the world's Muslim population, and included exceptions and waivers.

\textbf{4. Express Repudiation of \textit{Korematsu}.} Responding to Justice Sotomayor's dissent comparing the Proclamation to World War II Japanese-American internment, Roberts declared that \textit{Korematsu} ``has nothing to do with this case'' because forcibly relocating U.S. citizens to concentration camps solely by race is objectively unlawful. Seizing the occasion, the Court formally repudiated \textit{Korematsu} as having ``no place in law under the Constitution.''

\section{Concurring and Dissenting Opinions}

\begin{itemize}
  \item \textbf{Justice Kennedy (concurring):} Joined the majority in full but wrote separately to caution that judicial deference does not absolve executive officials of their independent constitutional duty to adhere to the First Amendment, even when their actions are ``not subject to judicial scrutiny or intervention.''
  \item \textbf{Justice Thomas (concurring):} Argued that the President also holds inherent Article II authority to exclude aliens and challenged the legality of ``universal injunctions'' (nationwide relief extending to non-parties), asserting that district courts lack historical equitable power under the Judiciary Act of 1789 to issue them.
  \item \textbf{Justice Breyer (dissenting, joined by Justice Kagan):} Focused on whether the Proclamation's exemption and waiver system was administered in good faith. Citing evidence that waivers were granted at tiny rates despite stated criteria, Breyer argued that if the waiver regime operated as ``window dressing,'' the Proclamation was tainted by religious animus and should be remanded for factual development while maintaining the injunction.
  \item \textbf{Justice Sotomayor (dissenting, joined by Justice Ginsburg):} Charged that a ``reasonable observer'' reviewing the President's repeated anti-Muslim statements would conclude the Proclamation's primary purpose was religious hostility, violating the Establishment Clause just as state commissioners' hostile remarks did in \textit{Masterpiece Cakeshop v. Colorado Civil Rights Commission} (2018). While welcoming the formal overruling of \textit{Korematsu}, Sotomayor warned that the majority ``redeploys the same dangerous logic'' by accepting a pretextual national-security rationale to validate group discrimination.
\end{itemize}

\section{Key Quotations}

\begin{quote}
``By its terms, \S~1182(f) exudes deference to the President in every clause. It entrusts to the President the decisions whether and when to suspend entry, whose entry to suspend, for how long, and on what conditions.''
--- Chief Justice Roberts, \href{https://supreme.justia.com/cases/federal/us/585/17-965/}{585 U.S. at 684}.
\end{quote}

\begin{quote}
``Plaintiffs challenge the entry suspension based on their perception of its effectiveness and wisdom, but the Court cannot substitute its own assessment for the Executive's predictive judgments on such matters.''
--- Chief Justice Roberts, \href{https://www.law.cornell.edu/supremecourt/text/17-965}{585 U.S. at 708}.
\end{quote}

\begin{quote}
``The forcible relocation of U.S. citizens to concentration camps, solely and explicitly on the basis of race, is objectively unlawful and outside the scope of Presidential authority. ... \textit{Korematsu} was gravely wrong the day it was decided, has been overruled in the court of history, and---to be clear---`has no place in law under the Constitution.'''
--- Chief Justice Roberts, \href{https://supreme.justia.com/cases/federal/us/585/17-965/}{585 U.S. at 710} (quoting \textit{Korematsu}, 323 U.S. at 248 (Jackson, J., dissenting)).
\end{quote}

\begin{quote}
``By blindly accepting the Government's misguided invitation to sanction a discriminatory policy motivated by animosity toward a disfavored group, all in the name of a superficial claim of national security, the Court redeploys the same dangerous logic underlying \textit{Korematsu} and merely replaces one `gravely wrong' decision with another.''
--- Justice Sotomayor (dissenting), \href{https://www.supremecourt.gov/opinions/17pdf/585us2r70_jgkn.pdf}{585 U.S. at 754}.
\end{quote}

\section{Doctrinal Significance and Modern Legacy (Through 2026)}

\textit{Trump v. Hawaii} cemented sweeping executive authority over foreign-national entry and became a pivotal precedent across multiple branches of separation-of-powers law through 2026:
\begin{itemize}
  \item \textbf{Executive Immigration Authority:} Although President Biden revoked Proclamation No. 9645 on January 20, 2021 (Proclamation No. 10141), \textit{Hawaii}'s expansive reading of \S~1182(f) remained controlling law and provided the legal template for renewed nationality-based entry restrictions issued during President Trump's second term in 2025--2026 (\href{https://www.americanimmigrationcouncil.org/report/trump-2025-travel-ban/}{American Immigration Council, Aug. 2025}).
  \item \textbf{Consular Nonreviewability and Presidential Speech:} In \href{https://supreme.justia.com/cases/federal/us/602/23-334/}{\textit{Department of State v. Mu\~noz}, 602 U.S. 899 (2024)}, the Court reinforced \textit{Hawaii} and \textit{Mandel} in holding that a U.S. citizen has no fundamental liberty interest in a noncitizen spouse's visa admission. That same month, in \href{https://www.law.cornell.edu/supremecourt/text/23-939}{\textit{Trump v. United States}, 603 U.S. 593 (2024)}, the Court quoted \textit{Hawaii} (585 U.S. at 701) for the President's ``extraordinary power to speak to his fellow citizens and on their behalf'' in defining official-act presidential immunity.
  \item \textbf{End of Universal Injunctions:} Justice Thomas's concurrence in \textit{Hawaii} (585 U.S. at 713--714) proved prophetic: on June 27, 2025, in \href{https://www.supremecourt.gov/opinions/24pdf/24a884_8n59.pdf}{\textit{Trump v. CASA, Inc.}, 606 U.S. 831 (2025)}, the Supreme Court adopted Thomas's historical equity framework and held that federal district courts generally lack authority under the Judiciary Act of 1789 to issue universal injunctions extending beyond the parties before them.
  \item \textbf{Anti-Canon and Rational-Basis Deference:} In \textit{Students for Fair Admissions v. Harvard} (2023), the Court confirmed that \textit{Hawaii} formally overruled \textit{Korematsu}. Meanwhile, through 2026, federal litigants continue invoking \textit{Hawaii}'s deferential rational-basis test in challenges to executive national-security and military-personnel policies.
\end{itemize}

\section{Pedagogical Critiques and Exam Hypotheticals}

\textbf{Pedagogical Critiques.}
\begin{itemize}
  \item \textbf{Double Standard on Religious Animus?} Critics highlight the jarring contrast between \textit{Masterpiece Cakeshop} (decided three weeks earlier in June 2018), where hostile remarks by state civil-rights commissioners invalidated a neutral order, and \textit{Trump v. Hawaii}, where far more explicit presidential statements were largely discounted under national-security deference.
  \item \textbf{Nondelegation Tension:} Reading \S~1182(f) as a nearly boundless grant to suspend entry of any class of aliens for any duration sits uneasily alongside the Court's modern ``major questions'' and nondelegation skepticism in domestic regulatory cases, underscoring a sharp foreign-affairs exceptionalism.
  \item \textbf{Symbolic Overruling of \textit{Korematsu}:} While formally burying \textit{Korematsu} removed a notorious stain from the U.S. Reports, scholars debate whether repudiating a 1944 domestic racial-internment precedent in dicta offers meaningful protection if courts defer uncritically to executive national-security assertions at the border.
\end{itemize}

\textbf{Exam Hypothetical 1 (Direct \textit{Trump v. Hawaii} Application).} Following a six-month interagency review finding that four foreign nations---three of which are predominantly Buddhist---lack biometric passport controls and fail to share criminal-record databases, the President issues a proclamation under 8 U.S.C. \S~1182(f) suspending nonimmigrant tourist visas from those four nations, with case-by-case hardship waivers. During the campaign, the President repeatedly vowed to ``halt Buddhist immigration.'' A U.S. university challenges the proclamation under the INA and Establishment Clause. \textit{Analysis:} Valid under \textit{Trump v. Hawaii}. Section 1182(f) confers broad discretion upon a presidential finding of detriment, and under \textit{Mandel}/rational-basis review, the facially neutral text, interagency vetting record, and waiver provisions supply an independent national-security justification sufficient to survive Establishment Clause scrutiny.

\textbf{Exam Hypothetical 2 (``Opposite Side'' Contrast---Domestic Application and \textit{Masterpiece Cakeshop}).} By contrast, suppose the same President issues an executive order directing the Department of Education to revoke federal accreditation eligibility from all \textit{domestic} private colleges located in the United States that offer Buddhist-studies degree programs, citing a vague ``national cultural security'' concern after making identical campaign statements denigrating Buddhism. \textit{Analysis:} Unconstitutional. Unlike \textit{Trump v. Hawaii}, this policy regulates domestic institutions and U.S. citizens within the United States rather than foreign-national entry at the border, and it lacks any statutory delegation akin to \S~1182(f). Without \textit{Mandel}'s border-entry deference, domestic Establishment Clause and Free Exercise scrutiny (\textit{Masterpiece Cakeshop}; \textit{Church of the Lukumi Babalu Aye}) applies in full force, and the overt religious targeting and extrinsic evidence of animus render the order invalid.

\end{document}
